\section{A Comparative Empirical Case Study}
\label{sec:exp}

In this section we describe an empirical case study in which we
apply our regularization framework to six different datasets
in which fairness is a central concern. These datasets include
cases in which the observed labels are real-valued, and cases in
which they are binary-valued. For the real-valued datasets, we apply
standard linear regression with our various fairness regularizers.
For the binary-valued datasets, we apply standard logistic regression,
again along with fairness regularizers. 
% This
% allows us to interpret our model predictions as probabilities of a 
% binary outcome (such as defaulting on a loan or committing a crime) via
% a proper scoring rule,
% which seems appropriate for these datasets (as opposed to 
% settings in which the ground truth is truly binary, e.g.
% being older or younger than 21 years old).
For datasets with real-valued targets we
normalized the inputs and outputs to be zero mean and unit
variance, and we set the cross-group fairness weights as
$d(\yi, \yj) = e^{-(\yi-\yj)^2}$;
for datasets with binary targets 
we set
$d(\yi,\yj) = \mathbbm{1}[\yi = \yj]$.

For each dataset $S$, our framework requires that we 
solve optimization problems of the
form $\min_{\vw} \ell(\vw,S) + \lambda f(\vw, S) + \gamma ||\vw||_2$ for
variable values of $\lambda$,
where $\ell(\vw,S)$ is either MSE (linear regression)
or the logistic regression loss.
For each $\lambda$ we picked $\gamma$ as
a function of this $\lambda$ by cross validation (see Appendix~\ref{sec:cross} for more details).
All optimization
problems are solved using the CVX solver in Matlab (for real-valued datasets) or
python (for binary-valued datasets).\footnote{See \url{http://www.cvxr.com}~and~\url{http://www.cvxpy.org} for more details. We set the number
of iterations to be 1000 in our optimization solvers.}
Furthermore, all the results are reported using 10-fold cross
validation (see more details in Appendix~\ref{sec:deets}).

The datasets themselves are summarized in Table~\ref{tab:summary},
where we specify the size and dimensionality of each, along with 
the ``protected'' feature (race or gender) 
that thus defines the subgroups across which we apply our
fairness criteria (see Appendix~\ref{sec:data-description} for more details). 
The datasets vary considerably in the
number of observations, their dimensionality, and the relative
size of the minority subgroup.

The \emph{Adult} dataset~\citep{uci-repo, adult-dataset} from the UC Irvine Repository contains 1994
Census data, and the goal is to predict whether the income of an
individual in the dataset is more than $50$K per year or not. 
The sensitive or protected attribute is gender.\footnote{We only used the data in Adult.data in our experiments.}
The \emph{Communities and
  Crime} dataset \citep{uci-repo} includes features relevant to
per capita violent crime rates in different communities in the United
States, and the goal is to predict this crime
rate; race is the protected variable.
The \emph{COMPAS}
dataset contains data from Broward County, Florida, originally compiled
by ProPublica~\cite{propublica}, in which the goal is to predict
whether a convicted individual would commit a violent crime in the
following two years or not. The protected attribute 
is race, and the data was filtered in a fashion similar to that of~\citet{Corbett-DaviesP17}.
The \emph{Default}
dataset~\citep{default-dataset,uci-repo} contains data from Taiwanese
credit card users, and the goal is to predict whether an individual
will default on payments. The protected attribute is
gender.
The \emph{Law School} dataset\footnote{\url{http://www2.law.ucla.edu/sander/Systemic/Data.htm}} 
consists of the records of law students who went on to take the bar exam. The goal is to 
predict whether a student will pass the exam based on features such as LSAT score and 
undergraduate GPA. The protected attribute is gender. 
The \emph{Sentencing} dataset contains information from a state department
of corrections regarding inmates in 2010. The goal is
to predict the sentence length given by the judge based on factors
such as previous criminal records and the crimes for which the
conviction was obtained. The protected attribute is gender.


\begin{table}[h]
\begin{center}
\begin{tabular}{|l|c|c|c|c|c|c|}
\hline
{\bf Data Set\/} & {\bf Type\/}  & {\bf $n$\/} & {\bf $d$\/} & {\bf Minority $n$\/} & {\bf Protected\/}
   \\ \hline
Adult & logit & 32561 & 14 & 10771 & gender
     \\\hline
Communities and Crime & linear & 1994 & 128 & 227 & race
     \\\hline
COMPAS & logit & 3373 & 19 & 1455 & race
     \\\hline
Default & logit & 30000 & 24 & 11888 & gender
     \\\hline
%     Student &  &  \checkmark &395 & & 187 %&7.17\% (20.69\%) 
 %    \\\hline
Law School & logit & 27478 & 36 & 12079 & gender
     \\\hline
%     Stop and Frisk & linn & 29786 & & 5467 %& 13.46\% (28.62\%) 
   %  \\\hline
Sentencing & linear & 5969 & 17  & 385 & gender
     \\\hline
\end{tabular}
\end{center}
\caption{Summary of datasets. 
Type indicates whether regression is logistic or linear;
$n$ is total number of data points;
$d$ is dimensionality;
Minority $n$ is the number of data points in the smaller population;
Protected indicates which feature is protected or fairness-sensitive.
}
\label{tab:summary}
\end{table}
